% Literature notes for the taxes_ads project.
% Raw material for the paper, not paper prose.

\section*{Related work}

\subsection*{\citet{keller_guyt_grewal_soda_marketing_conduct}
  --- Soda Taxes and Marketing Conduct, JMR 61(3), 2024}

\paragraph{Question.} After a soda tax, do retailers change
marketing conduct beyond price (promo frequency, promo depth,
feature frequency), and does consumer sensitivity to those
instruments change?

\paragraph{Data.} Nielsen RMS scanner data (Kilts), UPC--store--week.
Five treated localities (Philadelphia, Berkeley, Boulder,
Oakland, San Francisco --- inferred from the paper's market
list); 208 stores, 13{,}690 UPC--store pairs, 587 UPCs, 167
brands. Sample window: vote date through one year post-tax,
dropping the first three weeks post-introduction
(``dust-settling'' convention; Dekimpe et al.\ 2008).
UPC--store pairs kept only if consistently offered in both the
treated and matched-control markets and with at most 25\% of
weeks zero-sales. 96\% of taxed UPCs appear in both treated
and control markets, covering $>$99\% of volume.

\paragraph{Findings: conduct changes (treated vs.\ matched control).}
\begin{itemize}
  \item Promotion frequency: $-2\%$.
  \item Promotion depth: $-12\%$.
  \item Feature promotion frequency: $-14\%$.
  \item Regular price increases consistent with prior pass-through
        literature.
  \item Net effect: non-price conduct changes worsen the tax's
        negative sales effect rather than offsetting it.
\end{itemize}

\paragraph{Findings: effectiveness changes.}
Post-tax, consumers become \emph{less} sensitive to regular price,
feature promotions, and promo depth, but \emph{more} sensitive
to promotion incidence. Conduct response and effectiveness move
in opposite directions on the promotion-incidence margin
(consumers respond more to promos, managers run fewer).

\paragraph{Cross-store rigidity within market.}
Average correlation across stores within a UPC--retailer--market
of regular price (promo incidence, promo depth):
$0.265\ (0.164,\ 0.310)$ pre-tax,
$0.402\ (0.159,\ 0.348)$ post-tax. Some rigidity, but more
local adaptation than DellaVigna and Gentzkow (2019) or
Hitsch, Horta\c{c}su, and Lin (2021). Features are set at the DMA
level and so are constant across stores within a
UPC--retailer--week.

\paragraph{Event panel (Table 1 of the paper).}
A useful ready-made list of U.S.\ SSB tax events. Pass-through
column lists rates reported in the cited studies.

\begin{tabular}{lllrll}
\toprule
Locality & Vote & Implementation & Tax (\textcent/oz.) & Pass-through & Products \\
\midrule
Philadelphia & 2016-06-16 & 2017-01-01 & 1.5 &
  103\% (Cawley et al.\ 2020); 90\% (Roberto et al.\ 2019);
  97\% (Seiler, Tuchman, Yao 2021) &
  Sweetened beverages incl.\ regular and diet soda \\
Boulder & 2016-11-08 & 2017-07-01 & 2.0 &
  53--154\% (Cawley et al.\ 2021) &
  SSBs $\geq$ 5g caloric sweetener / 12 fl.\ oz. \\
Oakland & 2016-11-08 & 2017-07-01 & 1.0 &
  92\% (Falbe et al.\ 2020) &
  SSBs $\geq$ 25 cal / 12 fl.\ oz. \\
San Francisco & 2016-11-08 & 2018-01-01 & 1.0 &
  100\% (Falbe et al.\ 2020) &
  SSBs $>$ 25 cal / 12 fl.\ oz. \\
Seattle & 2017-06-05 & 2018-01-01 & 1.75 &
  90\% (Jones-Smith et al.\ 2020); 59\% (Powell and Leider 2020) &
  Sweetened beverages $\geq$ 40 cal / 12 fl.\ oz. \\
\bottomrule
\end{tabular}

Cook County / Chicago (2017 SSB tax, repealed) is not in this
panel and has the bottled-water-is-treated wrinkle noted in
\texttt{CLAUDE.md}.

\paragraph{Relevance to \texttt{taxes\_ads}.}
This paper is the closest precedent for ``does conduct beyond
price respond to a category tax.'' Two contrasts:
\begin{itemize}
  \item They use retailer marketing conduct (in-store promo,
        feature) measured at UPC--store--week. We measure
        \emph{manufacturer} advertising conduct (Nielsen Ad
        Intel Spot TV) at brand--DMA--month.
  \item They find non-price conduct \emph{shrinks} post-tax in
        the treated market. The analogous prediction for us is
        that national advertisers cut Spot TV buys in the
        treated DMA. The Philly extraction
        (\texttt{experiments/philly\_soda\_tax/}) is the first
        test of that.
\end{itemize}

\subsection*{\citet{hua_philly_beverage_ad_expenditures}
  --- Philadelphia Beverage Tax's Impact on Beverage Ad
  Expenditures and Number of Ads Purchased, AJPM, 2024}

\paragraph{Question.} Did beverage ad expenditures and number
of ads purchased change in Philadelphia versus Baltimore after
the Jan 1, 2017 Philadelphia beverage excise tax?

\paragraph{Data and design.} Monthly beverage ad expenditures
and ad counts by brand, Jan 2016 -- Dec 2019, coded as taxed
vs.\ not taxed (analyzed 2023). Controlled interrupted time
series on segmented linear regressions; outcomes are quarterly
taxed-beverage ad expenditures and number of ads, aggregated
and stratified by internet, spot TV, and local radio.

\paragraph{Findings.}
\begin{itemize}
  \item No significant differences in taxed-beverage ad
        expenditures between Philadelphia and Baltimore on
        pre-trend, at implementation, or post-trend.
  \item $0.13$ fewer quarterly taxed ads per 100 households in
        Philadelphia vs.\ Baltimore at baseline (95\% CI:
        $-0.25, -0.003$).
  \item Internet: $0.42$ fewer quarterly taxed ads per 100
        households in Philadelphia immediately post-tax (95\%
        CI: $-0.77, -0.06$).
  \item Spot TV: share of taxed-beverage ads was \emph{higher}
        in Philadelphia at baseline by 28.0 pp (95\% CI:
        $1.9, 54.1$).
\end{itemize}

\paragraph{Authors' conclusion (verbatim).}
``This study found little evidence of changes in mass media
advertising on the examined platforms between 2016 and 2019
due to the Philadelphia beverage tax.''

\paragraph{Relevance to \texttt{taxes\_ads}.}
\begin{itemize}
  \item Closest direct precedent: same event (Philly 2017),
        same outcome class (beverage advertising), Baltimore
        as the comparison DMA --- the user's own short list
        for our extraction includes Baltimore.
  \item Their null is a meaningful prior. If the AJPM result
        is right, the Spot TV margin we are extracting should
        also show a null in Philly. Our extraction is broader
        (national panel, all parents, multiple SSB taxes
        eventually), so even if Philly alone is null, pooled
        evidence may move.
  \item Methodological contrast: segmented ITS with one
        comparison city vs.\ our stacked DiD with multiple
        events and a synthetic / multi-DMA control. Different
        identification, different power.
  \item Spot TV baseline difference (Philly higher by 28 pp in
        taxed-beverage share) is worth replicating from the
        Ad Intel panel before running the event study; large
        baseline gaps cast doubt on the parallel-trends
        assumption.
\end{itemize}

\subsection*{\citet{dubois_griffith_oconnell_junk_food_ads}
  --- The Effects of Banning Advertising in Junk Food Markets,
  ReStud 85(1), 2018}

\paragraph{Question.} If advertising of junk food is banned,
does diet quality improve? Answer depends on how advertising
shifts demand and how firms reprice in response.

\paragraph{Setting and data.} U.K.\ potato chips market.
Consumer-level exposure to TV adverts merged with household
purchase data. Demand allows advertising to shift the weight
on product healthiness, tilt demand curves, generate dynamics,
and spill over across brands.

\paragraph{Findings.} A ban's health benefit is \emph{partially
offset} on two margins: (i) firms cut prices in response to the
demand shift, and (ii) consumers substitute toward other junk
foods.

\paragraph{Relevance to \texttt{taxes\_ads}.}
\begin{itemize}
  \item Direct evidence that firms reprice when their
        advertising channel is shut down. The mirror prediction
        for a tax: if the tax suppresses advertising (our
        outcome), prices may move in ways not captured by
        pass-through studies that hold advertising fixed.
  \item Cross-brand and cross-category substitution matter for
        welfare. Our event panel only measures the
        advertising-response margin; combining with scanner
        pass-through and substitution estimates would close the
        loop.
  \item Methodologically distant (structural demand, ad-exposure
        micro-data) but the question --- how does the
        advertising margin interact with category policy --- is
        the same.
\end{itemize}
